\chapter*{Введение}
\addcontentsline{toc}{chapter}{Введение}
\label{ch:intro}

Робототехнические системы, предназначенные для манипулирования объектами, занимают центральное место в современных исследованиях на стыке робототехники и искусственного интеллекта. Задачи захвата, перемещения и точного позиционирования предметов возникают в широком спектре приложений: от промышленной сборки и складской логистики до бытового обслуживания и медицинской хирургии. На протяжении последних двух десятилетий подходы к решению этих задач эволюционировали от жестко запрограммированных контроллеров, опирающихся на точные модели среды, к методам, способным адаптироваться к изменяющимся условиям благодаря обучению на данных.

Одним из наиболее перспективных направлений является визуомоторное управление -- подход, при котором политика манипулятора формируется непосредственно на основе визуальных наблюдений с камер. Парадигма имитационного обучения, в частности клонирование поведения (behavior cloning), позволяет обучать такие политики на основе демонстраций оператора-человека, без необходимости проектировать функцию вознаграждения и длительно взаимодействовать с симулятором. За последние годы в этой области были предложены архитектуры, существенно продвинувшие качество обученных политик: рекуррентные сети, диффузионные политики, трансформерные модели поведения.

Тем не менее практическое применение визуомоторных политик сопряжено с серьезными трудностями. Действия манипулятора представляют собой многомерные непрерывные векторы с потенциально мультимодальным распределением. Ошибки предсказания накапливаются во времени, приводя к расхождению траекторий. Особенно критична проблема визуальных доменных сдвигов: изменения освещения, фона, появление посторонних объектов или частичные окклюзии могут существенно исказить входные представления и привести к отказу обученной политики. Современные архитектуры используют высокоразмерные непрерывные представления изображений, которые, как правило, содержат избыточную информацию, чувствительную к подобным возмущениям.

В этой связи возникает вопрос: возможно ли повысить устойчивость визуомоторных политик путем структурирования промежуточных представлений? Одним из инструментов такого структурирования является векторное квантование -- метод, отображающий непрерывные векторы на конечный набор дискретных прототипов. В последние годы квантование представлений продемонстрировало свою эффективность в ряде робототехнических приложений -- для дискретизации действий, состояний и визуальных наблюдений. Однако систематическое исследование влияния квантования именно визуальных представлений на обучаемость (эффективность при малом числе демонстраций) и надежность (устойчивость к визуальным доменным сдвигам) алгоритмов захвата объектов до настоящего времени не проводилось.

Целью настоящей работы является экспериментальная оценка влияния квантования векторных представлений на качество обучения и устойчивость алгоритмов захвата объектов манипулятором с визуальной обратной связью.

Работа организована следующим образом. В первой главе приведен аналитический обзор литературных источников, охватывающий основные подходы к имитационному обучению для задач манипулирования, методы квантования представлений и их приложения в робототехнике. Во второй главе описана методология исследования, включающая выбор базового алгоритма, архитектуру модуля квантования и методику проведения экспериментов. Третья глава посвящена результатам вычислительных экспериментов и их анализу. В заключении сформулированы основные выводы и рекомендации.

\endinput
